\todoDraftOnly{Revisit counting work section in a bit and mb redraft}

\todoDraftOnly{review all of 3.1 for references to PoRs only providing weight to the best block, etc. (this was an old idea and is not the case.)}

\providecommand\BL[1]{\BCI{L}{i}{#1}}
\providecommand\BR[1]{\BCI{R}{j}{#1}}
% note: \BR renewed below to change `j` term to `k`

When chains L and R have the same block frequency and produce blocks in an orderly fashion, it's trivial to count how much work is contributed by each PoR.
However, real-world blockchains do not produce blocks like this --- even if chains L and R have the same block frequencies, sometimes L will produce 2 or 3 (or more) blocks before R produces its next one (or vice versa).
\begin{equation*}
    \xyMChains{
        \xyL{\text{Chain L}} & \cdots & \xyB{L_i} \ar[l] & & \xyB{L_{i+1}} \ar[ll] \ar@{.>}[ld] & \xyB{L_{i+2}} \ar[l] & \xyB{L_{i+3}} \ar[l] & \\
        \xyL{\text{Chain R}} & \cdots & & \xyB{R_k} \ar[ll] \ar@{.>}[lu] & & & & \xyB{R_{k+1}} \ar[llll] \ar@{.>}[lu] \\
        \xyTimeHoriz{rrrrrrr} & & & & & & & \\
    }
\end{equation*}
Should \BR{1} include 3 PoRs, or just 1?
How should \BR{1} calculate the weight of those PoRs and the total weight contributed to chain R?

If \BR{1} had reflected only \BL{1}, then it's clear that we would count work from both \BL{1} and \BR{1} --- that much is easy.
However, \BR{1} reflected \emph{three} L headers.
Should \BR{1} count work from \emph{all three}, or something else?
Why?

Consider the situation from the miner's perspective (the miner of \BL{2}).
There is no new R block for them to reflect --- if there were, they'd include it.
The miner is effectively claiming --- via the absence of a new reflection --- that \BL{2}'s parent (\BL{1}) reflects the most recent R header ($R_k$); that \BL{1} is still \emph{up-to-date}.
From this perspective, it's clear that \BL{2} should contribute reflected weight to the best R headers that were most recently reflected: in this case, just $R_k$.

We can evaluate \BL{3} the same way.

So, yes \BR{1} should include and count work from all 3 PoRs in this case.

\providecommand\BRj[1]{\BCI{R}{j}{#1}}
\renewcommand\BR[1]{\BCI{R}{k}{#1}}

What about a more complicated example?
\begin{equation*}
    \xyMChains{
        \xyL{\text{Chain L}} & \cdots & \xyB{L_i} \ar[l] & & \xyB{L_{i+1}} \ar[ll] \ar@{.>}[ld] & \xyB{L_{i+2}} \ar[l] & \xyB{L_{i+3}} \ar[l] & & \xyB{L_{i+4}} \ar[ll] \ar@{.>}[ld] \\
        \xyL{\text{Chain R}} & \cdots & & \xyB{R_k} \ar[ll] \ar@{.>}[lu] & & \xyB{R_{k+1}} \ar[ll] & & \xyB{R_{k+2}} \ar[ll] \ar@{.>}[lu] & \\
        \xyTimeHoriz{rrrrrrrr} & & & & & & & & \\
    }
\end{equation*}
\BR{2} can be evaluated by similar logic to \BRj{1} above, but what about \BL{4}?

Does \BL{4} count the weight of \BR{1}?
If so, does it count towards \BL{3}'s chain-weight?

The difficulty in answering these questions --- with what we've covered so far --- is that \emph{the fork-rule doesn't care} where the weight is added, just that it \emph{is} added.
When the fork rule compares \BL{4} to an alternative block, only the \emph{total} chain-weight of each is relevant.
However, to understand exactly how PoRs \emph{should} be counted, we need to know more.

The essence of PoR is \emph{the confirmation of another chain's \ul{history}} --- so a reflection cannot contribute to a block which exists \emph{in the reflecting block's \ul{future}}.
It must be that reflections contribute to the \emph{most recent common ancestor} of all \emph{possible} L blocks which \emph{could} include that specific PoR.

This leads us to a simple and elegant way to count reflected work: follow the arrows back from the reflecting block to the most recently reflected local block.

\BL{4}'s reflection of \BR{2} contributes chain-weight to \BL{3} because \emph{any} block that builds on \BL{3} could reflect \BR{2}.
% (L_{i+4} \; \ar@{.>}[r] &) \;
Following the arrows: $\xyInlineRefl{R_{k+2} \; \ar@{.>}[r] & \; L_{i+3}}$.
% Note that the blocks and links in parentheses are specific to the block

% Following the arrows: $L_{i+4} \dashrightarrow R_{k+2} \dashrightarrow L_{i+3}$.

Similarly, \BL{4}'s reflection of \BR{1} (which could be either explicit or implicit via \BR{2}) must contribute chain-weight to $L_i$ because \emph{any} block building on $L_i$ could reflect \BR{1} and include that PoR.
% (L_{i+4} \; \ar@{.>}[r] & \; R_{k+2} \; \ar[r] &) \;
Following the arrows: $\xyInlineRefl{R_{k+1} \; \ar[r] & \; R_k \; \ar@{.>}[r] & \; L_{i}}$.

% \BL{4} \textrightarrow \BR{2} \textrightarrow \BR{1} \textrightarrow $R_k$ \textrightarrow $L_i$.
% Following the arrows: $L_{i+4} \dashrightarrow R_{k+2} \to R_{k+1} \to R_k \dashrightarrow L_{i}$.

We can also say more about \BR{2}'s reflection of \BL{3}, now, too.
Particularly, that the reflected work of $\BL{1}\cdots \BL{3}$ should contribute to \BR{}, \emph{not} \BR{1}.

\begin{draftonly}

\todoDraftOnly{subsubsec counting work extra}

\todoDraftOnly{
consider a situation where an attacker reflects 1 honest block and 1 competing attacker block, and the attacker block is not released.
for the sake of the argument, assume the reflected blocks weight much less than the reflecting block.
}


\todoDraftOnly{
    consensus enforces availability of refl blocks
    but counting work doesn't know anything about 'availability'
    so the only thing it can do is know if something is within its history or not.
    asymmetry.
}

\todo{
    noncontributing pors: like ghost (uncles) for PoR --- reflection isn't 2 ways?
    is this okay with miner incentives? or is there a conflict?
}

\end{draftonly}
